\documentclass[]{article}
\usepackage{lmodern}
\usepackage{amssymb,amsmath}
\usepackage{ifxetex,ifluatex}
\usepackage{fixltx2e} % provides \textsubscript
\ifnum 0\ifxetex 1\fi\ifluatex 1\fi=0 % if pdftex
  \usepackage[T1]{fontenc}
  \usepackage[utf8]{inputenc}
\else % if luatex or xelatex
  \ifxetex
    \usepackage{mathspec}
  \else
    \usepackage{fontspec}
  \fi
  \defaultfontfeatures{Ligatures=TeX,Scale=MatchLowercase}
\fi
% use upquote if available, for straight quotes in verbatim environments
\IfFileExists{upquote.sty}{\usepackage{upquote}}{}
% use microtype if available
\IfFileExists{microtype.sty}{%
\usepackage{microtype}
\UseMicrotypeSet[protrusion]{basicmath} % disable protrusion for tt fonts
}{}
\usepackage[margin=1in]{geometry}
\usepackage{hyperref}
\hypersetup{unicode=true,
            pdftitle={NHTSA DSFS Meta-Analysis},
            pdfborder={0 0 0},
            breaklinks=true}
\urlstyle{same}  % don't use monospace font for urls
\usepackage{longtable,booktabs}
\usepackage{graphicx,grffile}
\makeatletter
\def\maxwidth{\ifdim\Gin@nat@width>\linewidth\linewidth\else\Gin@nat@width\fi}
\def\maxheight{\ifdim\Gin@nat@height>\textheight\textheight\else\Gin@nat@height\fi}
\makeatother
% Scale images if necessary, so that they will not overflow the page
% margins by default, and it is still possible to overwrite the defaults
% using explicit options in \includegraphics[width, height, ...]{}
\setkeys{Gin}{width=\maxwidth,height=\maxheight,keepaspectratio}
\IfFileExists{parskip.sty}{%
\usepackage{parskip}
}{% else
\setlength{\parindent}{0pt}
\setlength{\parskip}{6pt plus 2pt minus 1pt}
}
\setlength{\emergencystretch}{3em}  % prevent overfull lines
\providecommand{\tightlist}{%
  \setlength{\itemsep}{0pt}\setlength{\parskip}{0pt}}
\setcounter{secnumdepth}{0}
% Redefines (sub)paragraphs to behave more like sections
\ifx\paragraph\undefined\else
\let\oldparagraph\paragraph
\renewcommand{\paragraph}[1]{\oldparagraph{#1}\mbox{}}
\fi
\ifx\subparagraph\undefined\else
\let\oldsubparagraph\subparagraph
\renewcommand{\subparagraph}[1]{\oldsubparagraph{#1}\mbox{}}
\fi

%%% Use protect on footnotes to avoid problems with footnotes in titles
\let\rmarkdownfootnote\footnote%
\def\footnote{\protect\rmarkdownfootnote}

%%% Change title format to be more compact
\usepackage{titling}

% Create subtitle command for use in maketitle
\newcommand{\subtitle}[1]{
  \posttitle{
    \begin{center}\large#1\end{center}
    }
}

\setlength{\droptitle}{-2em}
  \title{NHTSA DSFS Meta-Analysis}
  \pretitle{\vspace{\droptitle}\centering\huge}
  \posttitle{\par}
  \author{}
  \preauthor{}\postauthor{}
  \date{}
  \predate{}\postdate{}


\begin{document}
\maketitle

\centering
Meta-Analysis of Effectiveness of Dynamic Speed Feedback Signs:

Daniel Flynn Andrew Breck Olivia Gillham Donald L. Fisher

April 2017

Volpe National Transportation Systems Center 55 Broadway Cambridge, MA
02114

\clearpage
\tableofcontents
\clearpage
\raggedright

\section{Introduction}\label{introduction}

Traffic safety engineers and transportation planners have a high
interest in the efficacy of different approaches to reducing driver
speeds on critical roadway segments. One of these approaches, in
particular monitoring traffic speeds in conjunction with installation of
Dynamic Speed Feedback Signs (DSFS), has recently been targeted by the
Governors' Highway Safety Association (GHSA). An ongoing review by the
Volpe National Transportation Systems Center (Volpe) of the studies
conducted to investigate DSFS indicates that the studies have employed
similar enough designs and are sufficiently numerous that a
meta-analysis to quantitatively summarize the results can be carried
out.

Meta-analysis is a statistical technique used to draw general
conclusions from a set of related studies focused on similar
experimental treatments and responses (Glass 1976). Given sufficiently
similar study designs, the size of the treatment effects can be pooled
to create a meta-data set, which is then analyzed with appropriate
weighting based on the variability and size of each study. This
meta-analysis aims to provide decision-makers with a quantitative
summary of how much vehicle speed is reduced following installation of a
DSFS. Factors such as location, timing, vehicle type, posted speed,
safety focus, and others are all considered.

\section{Methods}\label{methods}

\subsection{Data collection}\label{data-collection}

Candidate publications were identified by a search of the following
databases: Transport Research International Documentation (TRID),
National Transportation Library (NTL), WorldShare, Academic Search
Complete, PsycINFO, Web of Science, and Science Direct. The searches
were carried out between 3/9/2016 and 3/16/2016 by research librarians.
A full list of the search terms is provided in Table 1. Data were
entered at two levels:

\begin{enumerate}
\def\labelenumi{(\alph{enumi})}
\tightlist
\item
  the descriptor level, which describes the characteristics of each
  study (the features which identify how the DSFS was implemented, the
  hypotheses that were evaluated, and the experimental design)
\item
  the data level, which was the least-aggregated level of data reported
  either in tabular or graphical format in the publication. For
  calculations of effect size, at a minimum a study would have to report
  the mean and standard deviation (SD) of vehicle speeds, as well as the
  number of vehicles observed, at either two points in time or two
  locations. The difference in whether studies focused on different
  timing or location of vehicle speed data collection determines which
  hypotheses can be tested (see section titled Calculating effect size).
\end{enumerate}

Publications in some cases included more than one study, meaning they
tested DSFS at different locations or in combinations with additional
treatments such as police cruisers or speed radar enforcement; these are
treated as separate studies in the data collection, grouped under the
same publication ID. Publications employed many different combinations
of descriptors in the study design, and effort was made to enter data
from any relevant descriptor that could be expected to appear in
multiple studies. Key descriptors included posted speed, safety focus,
vehicle type, sign type, and distance from sign the vehicle speeds were
measured.

\subsection{Calculating Effect Size}\label{calculating-effect-size}

Meta-analyses were initially used primarily to assess the odds of a
particular outcome, such as the odds a drug treatment would result in a
favorable outcome for a patient, and then were expanded to more
generally test the size and direction of the outcome across different
studies. Both types of tests, assessing odds of a categorical outcome or
size of a continuous outcome, are dependent on calculating the effect
size. For the meta-analysis of DSFS, the outcomes are vehicle speeds at
different points in time. Three sets of hypotheses regarding the
effectiveness of DSFS can tested, and each requires a different
calculation of effect size:

\begin{enumerate}
\def\labelenumi{\arabic{enumi}.}
\tightlist
\item
  H1. Activation of the DSFS. How much the installation of DSFS reduces
  driver speed is the central question for this analysis. The studies
  tested activation in three ways:

  \begin{enumerate}
  \def\labelenumii{\alph{enumii}.}
  \tightlist
  \item
    Compare vehicle speeds adjacent to the DSFS before activation and
    then during activation
  \item
    compare vehicle speeds upstream (of the DSFS) and adjacent (to the
    DSFS)
  \item
    Compare the difference between the upstream and adjacent vehicle
    speeds before activation to the difference between the upstream and
    adjacent vehicle speeds during activation of the DSFS. This is the
    ideal comparison, or the true effect (Cruzado and Donnell 2009).
  \end{enumerate}
\item
  H2. Activation of the DSFS: downstream. The studies tested the
  activation effect downstream in three ways:

  \begin{enumerate}
  \def\labelenumii{\alph{enumii}.}
  \tightlist
  \item
    Compare downstream vehicle speeds during activation to downstream
    vehicle speeds before activation,
  \item
    Compare downstream vehicle speeds adjacent to the DSFS during
    activation,
  \item
    Compare a difference of differences which controls for variations in
    both the time at which the vehicle speeds are measured and the
    location at which the vehicle speeds are measured.
  \end{enumerate}
\item
  H3. Deactivation of the DSFS. The comparisons are similar to the
  above, but now vehicle speeds are compared at the DSFS location after
  the DSFS has been removed to vehicle speeds at the DSFS while the DSFS
  was activated, or a normalized difference from after to before or
  after to during installation.
\end{enumerate}

We estimate the effect size (ES) associated with activation of the DSFS
appropriately for each hypothesis. The calculation of effect sizes are
based on the difference between the mean speeds adjacent to the DSFS
during the period in which it is active, \[\bar{X}_{DSFS}\] The sample
mean speed adjacent the DSFS before activation is defined as
\[ \bar{X}_{Before} \] The ES is this difference, placed relative to
units of the pooled standard deviation:
\[ ES = \frac{\bar{X}_{DSFS} - \bar{X}_{Before}}{SD_{Pooled}} \]

Where the pooled standard deviation for this calculation is defined as
follows:

\[
SD_{Pooled} = \sqrt{\frac{(N_{DSFS}-1)SD_{DSFS}}{N_{DSFS} + N_{Before} -2 } + \frac{(N_{Before}-1)SD_{Before}}{N_{DSFS} + N_{Before} -2 }}
\]

This effect size calculation is also referred to as the standardized
mean difference (Hedges and Olkin 1985). The pooled standard deviation
is calculated following standard calculations.

For other hypotheses based on the comparison of two groups, such as
comparing the effect of DSFS at two locations, adjacent to the DSFS or
upstream of the DSFS (Hypothesis H1C, Table 2), the effect size
calculation is similar.

Effect sizes based on normalized ``true effects'' require a slightly
modified calculation, where differences in means rather than means
themselves are compared. For example, when the activation effect of the
DSFS on vehicle speed is normalized for differences during and before
installation, at both the site adjacent to the sign and upstream of the
sign, the effect size is calculated as follows:

\[ ES = \frac{(\bar{X}_{Adjacent DSFS} - \bar{X}_{Adjacent Before}) - (\bar{X}_{Upstream DSFS} - \bar{X}_{Upstream Before})}{SD_{Pooled}} \]

The pooled standard deviation in this case is computed as follows:

\[
SD_{Pooled} = \sqrt{\frac{(N_{1}-1)SD_{1} + (N_{2}-1)SD_{2} + (N_{3}-1)SD_{3} + (N_{4}-1)SD_{4} }{N_{1}+N_{2} + N_{3} + N_{4} - 4}}
\]

The subscripts on each variable identify there relevant time and
location at the which the variable is measured: 1. Adjacent to the DSFS;
2. Adjacent, Before installation; 3. Upstream of the DSFS; and 4.
Upstream of the DSFS, Before installation. Figures 4-7 in the Literature
Review provide the detailed speed reduction formulas for each of the 13
possible combinations of timing and location of vehicle speeds in
response to DSFS installation. The specific calculations follow the two
forms above, for either the direct or normalized measurement of effect
size.

All of the effect size calculations and calculations of sample
variability \(v_{i}\) (see Hypothesis testing) were carried out in the
meta-analysis package \emph{metafor} (Viechtbauer 2010) in the
statistical programming language \emph{R} (R Core Team 2016). Hypothesis
testing (below) was carried out using the linear mixed-effect model
package \emph{lme4} (Bates, et al. 2015). The effect size calculations
in metafor have been validated against results provided by the software
packages Stata, SAS, and SPSS; the results from these other packages
either agreed completely or fell within a margin of error expected when
using numerical methods. The results of \emph{metafor} effect size
calculations and analyses have also been validated against textbook
examples.

\subsection{Hypothesis testing}\label{hypothesis-testing}

This meta-analysis investigates a single continuous outcome (speed of
vehicles) across multiple studies, with several moderating variables
(e.g., posted speed, safety focus, time of day). The use of multiple
moderating variables in a meta-analysis has also been called
meta-regression (van Houwelingen, Arends and Stijnen 2002). This type of
analysis is best conducted using a mixed-effect model, where studies are
the random effects, and the moderating variables are the fixed effects.
``Random effects'' refers to variables in the model which are assumed to
be representatives of a larger population (Raudenbush 2009, Hedges and
Vevea 1998), which in this case are the publications and studies within
each publication. The studies examined in this analysis are assumed to
be representatives of other possible studies that might have been done.
Statistically, random effects are variables which are themselves
modeled, with their own error terms. ``Fixed effects'' refers to
predictor variables in the model which are directly measured for each
study.

\[ y_{i} = \beta_{0} + \beta_{1}posted speed + \beta_{2} vehicle type+\beta_{3}safety focus + u_{i} + e_{i} \]

where \emph{i} indexes the study. \(\beta_{0}\) refers to the overall
intercept, while the terms \(\beta_{1}\), \(\beta_{2}\), and
\(\beta_{3}\) refer to the coefficients for each of teh predictor
variables.

The error term has two parts, first the study-level error term
\(u_{i}\), which is defined as \(u_{i} \sim N(0, \tau^{2})\), a normally
distributed error term with a mean of zero and variance \(\tau^{2}\);
this variance is estimated in the model, and is the random effect in the
model. As a random effect, it represents the variability assumed from
the larger population of studies from which the specific studies used
have been drawn. The remaining error term \(e_{i}\) is defined as
\(e_{i} \sim N(0, v_{i})\), meaning that the residual error is normally
distributed with a mean of zero and a variance of \(v_{i}\); this
variance \(v_{i}\) is directly computed in the effect size calculation.
The sampling variance \(v_{i}\) becomes important in the model fitting
process, giving greater weight to studies which have a lower variability
in the response, due to a large sample size, a low standard deviation in
the speeds, or a combination of the two.

For each of the 13 hypotheses tested, a similar structure of models was
used. Establishing which moderator variables (descriptors) should be
used will be the critical step in the analysis, as not all studies
provide complete information on all of the moderating variables.

This analysis allows results to be presented on the overall
effectiveness of DSFS in reducing speeds for a particular hypothesis,
for example the overall effect of DSFS activation when comparing speeds
measured at the DSFS sign to speeds measured at the site of the DSFS
sign, prior to sign installation (Hypothesis H1B). In addition to this
overall measure of speed reduction, this analysis will further allow
reporting of how much the posted speed drives the speed reduction, how
much the speed reduction differs for passenger cars versus commercial
vehicles, and how much the speed reduction differs for school zones,
work zones, rural-to-urban transition zones, or other safety focus
areas. The combination of multiple hypotheses tested and multiple
moderating variables should assist states in determining where and when
to deploy DSFS for maximum effect, and the size of the speed reduction
they can expect to observe.

\section{Results}\label{results}

\subsection{Summary of data}\label{summary-of-data}

The literature review identified 43 publications with data appropriate
for quantiative meta-analysis. From these publications, 57 studies were
identified, with a total of 204 sites at which the effectiveness of DSFS
were evaluated (Table \ref{tab:tab1}). By assessing availability of the
appropriate data (mean, \emph{s.d.}, and \emph{N}) within sites at the
appropriate locations, the number of publications, studies, and sites
available for testing each of hypothesis can be calculated. This shows
that hypothesis H1, the effect of DSFS activation at the site of the
sign, is the most commonly studied design, followed by H2, the effect of
DSFS activation downstream of the sign. Hypothesis H3, the effect of
DSFS following deactivation of the sign, was rarely studied, with only
three publications providing sufficient data for meta-analysis. Note
that these values differ from the Literature Review vote-count analysis,
which only requried studies report the significance of the result, not
necessarily providing the underlying data for a meta-analysis.

\begin{longtable}[]{@{}lrrr@{}}
\caption{\label{tab:tab1}Summary of the number of publications, studies,
and individual sites which can be used to test each of the 13
hypotheses.}\tabularnewline
\toprule
& Publications & Studies & Sites\tabularnewline
\midrule
\endfirsthead
\toprule
& Publications & Studies & Sites\tabularnewline
\midrule
\endhead
H1A & 12 & 22 & 81\tabularnewline
H1B & 18 & 36 & 186\tabularnewline
H1C & 14 & 28 & 123\tabularnewline
H2A & 9 & 18 & 85\tabularnewline
H2B & 9 & 18 & 85\tabularnewline
H2C & 9 & 18 & 85\tabularnewline
H2A' & 9 & 18 & 85\tabularnewline
H2C' & 9 & 18 & 85\tabularnewline
H3A & 2 & 3 & 8\tabularnewline
H3B & 3 & 5 & 8\tabularnewline
H3C & 2 & 3 & 4\tabularnewline
H3A' & 2 & 3 & 4\tabularnewline
H3B' & 3 & 5 & 8\tabularnewline
\bottomrule
\end{longtable}

\subsection{Activation hypothesis: H1}\label{activation-hypothesis-h1}

The three versions of hypothesis H1 all rely on a measurement of vehicle
speeds at the site of the DSFS, but use different baseline comparisons.
The simplest, H1B, compares vehicle speeds at the site of the sign
between the time before activation and during activation. H1A calculates
a similar value, for the difference between activation timing (before
and during), normalized to the same difference upstream of the DSFS. H1C
compares speeds at the DSFS to speeds upstream, at the same time.

\subsubsection{H1A: Activation
normalized}\label{h1a-activation-normalized}

Examining the mean effect sizes across vehicle types, without carrying
out any statiticaly modeling, it is clear that the largest decreases in
speed were observed for passenger cars.

\begin{longtable}[]{@{}lrrr@{}}
\caption{Decreases in speed for Hypothesis H1A, normalized activation
effect of DSFS, by vehicle type.}\tabularnewline
\toprule
Vehicle Type & Effect size yi & Variance vi & N\tabularnewline
\midrule
\endfirsthead
\toprule
Vehicle Type & Effect size yi & Variance vi & N\tabularnewline
\midrule
\endhead
Unspecified & -2.729 & 0.278 & 31\tabularnewline
Passenger & -3.905 & 1.035 & 16\tabularnewline
Truck & -1.973 & 4.401 & 25\tabularnewline
\bottomrule
\end{longtable}

\begin{longtable}[]{@{}lrrr@{}}
\caption{Summary of mixed effect model of hypothesis H1A, difference in
speeds betwen DSFS before and at time of activation, normalized to the
difference before and at time of activation upstream of the
sign.}\tabularnewline
\toprule
& Estimate & Std. Error & t value\tabularnewline
\midrule
\endfirsthead
\toprule
& Estimate & Std. Error & t value\tabularnewline
\midrule
\endhead
(Intercept) & -2.703 & 0.817 & -3.309\tabularnewline
Overall, the d & ecrease in & speed attribu & ted to DSFS, normalized by
upstream location, is -5.7 mph. Of the different grouping variables,
vehicle type has the greatest influence, with cars showing an addtional
-0.25 mph decrease. Trucks showed a smaller effect, with the decrease
being smaller by 0.30 mph. Only 28 sites were available for this
test.\tabularnewline
\bottomrule
\end{longtable}

\subsection{H1B}\label{h1b}

Calculations of effect sizes for H1B, the decrease in speed at the site
of the DSFS, show decreasese in speed across all vehicle types and
locations. The largest raw effect size was observed for passenger cars,
where speeds were on average reduced by -4.52 mph.

\begin{longtable}[]{@{}llrrr@{}}
\caption{Decreases in speed for Hypothesis H1B, activation effect of
DSFS, by location and vehicle type.}\tabularnewline
\toprule
Location & Vehicle Type & Effect size yi & Variance vi &
N\tabularnewline
\midrule
\endfirsthead
\toprule
Location & Vehicle Type & Effect size yi & Variance vi &
N\tabularnewline
\midrule
\endhead
1 upstream & Unspecified & -0.232 & 0.153 & 52\tabularnewline
2 adjacent & Unspecified & -2.687 & 0.136 & 130\tabularnewline
3 downstream & Unspecified & -3.072 & 0.195 & 28\tabularnewline
1 upstream & Passenger & -0.476 & 0.470 & 59\tabularnewline
2 adjacent & Passenger & -4.521 & 0.601 & 97\tabularnewline
3 downstream & Passenger & -2.725 & 0.411 & 72\tabularnewline
1 upstream & Truck & 0.394 & 2.167 & 92\tabularnewline
2 adjacent & Truck & -2.129 & 1.717 & 121\tabularnewline
3 downstream & Truck & -2.207 & 0.879 & 74\tabularnewline
\bottomrule
\end{longtable}

Statistical test of H1B, which takes into account the variation at each
level of site, study, and publication, as well as different vehicle
types and safety foci, shows a significant decrease in speeds of -3.42
mph at the location adjacent to the sign. Within passenger cars, this
decrease is even more pronounced with an estiamted additional decrease
of -0.55 mph. Thus, approximately 4 mph decrease is expected for
passenger cars at DSFS locations.

The reductions in speeds for trucks is more modest, with the speed
reduction smaller by 0.42 mph, so 3.00 mph decrease is expected for
trucks adjacent to DSFS.

\begin{longtable}[]{@{}lrrr@{}}
\caption{Summary of mixed effect model of hypothesis H1B, difference in
speeds betwen DSFS before and at time of activation.}\tabularnewline
\toprule
& Estimate & Std. Error & t value\tabularnewline
\midrule
\endfirsthead
\toprule
& Estimate & Std. Error & t value\tabularnewline
\midrule
\endhead
(Intercept) & -0.578 & 0.562 & -1.030\tabularnewline
location2 adjacent & -2.844 & 0.399 & -7.128\tabularnewline
\bottomrule
\end{longtable}

\begin{verbatim}
## Computing p-values via Kenward-Roger approximation. Use `p.kr = FALSE` if computation takes too long.
\end{verbatim}

\includegraphics{MA_Report_files/figure-latex/unnamed-chunk-4-1.pdf}

\subsection{H1C}\label{h1c}

The difference between speeds upstream and adjacent to DSFS were
surprisingly large, with large decreases in speed detected in several
studies which did not differentiate by vehicle type (Hallmark et al.
2007, Reddy et al. 2008, and Roberts et al. 2012). The mean
upstream-adjacent effect size for DSFS activation across all studies was
-2.73 mph.

\begin{longtable}[]{@{}lrrr@{}}
\caption{Decreases in speed for Hypothesis H1C, activation effect of
DSFS downstream, by vehicle type.}\tabularnewline
\toprule
Vehicle Type & Effect size yi & Variance vi & N\tabularnewline
\midrule
\endfirsthead
\toprule
Vehicle Type & Effect size yi & Variance vi & N\tabularnewline
\midrule
\endhead
Unspecified & -11.128 & 0.070 & 31\tabularnewline
Passenger & -1.704 & 1.585 & 16\tabularnewline
Truck & -1.424 & 2.920 & 25\tabularnewline
\bottomrule
\end{longtable}

\begin{longtable}[]{@{}lrrr@{}}
\caption{Summary of mixed effect model of hypothesis H1C, difference in
speeds between DSFS upstream and adjacent to the sign.}\tabularnewline
\toprule
& Estimate & Std. Error & t value\tabularnewline
\midrule
\endfirsthead
\toprule
& Estimate & Std. Error & t value\tabularnewline
\midrule
\endhead
(Intercept) & -8.615 & 3.254 & -2.647\tabularnewline
\bottomrule
\end{longtable}

\begin{verbatim}
## Computing p-values via Kenward-Roger approximation. Use `p.kr = FALSE` if computation takes too long.
\end{verbatim}

\includegraphics{MA_Report_files/figure-latex/unnamed-chunk-6-1.pdf}


\end{document}
